% !TeX root = ../../../main.tex

\subsubsection{Textual Time-Series Representation}\label{ch05:sec:textual-time-series}

{\raggedright
Given a clinical text document $T$, we extract a textual time series $S = \{(e_1,t_1), (e_2,t_2),...,(e_n,t_n)\}$ where $e_i$ represents a clinical finding (contiguous text span describing a temporally localizable patient event) and $t_i \in \mathbb{R}$ represents hours relative to a case-specific reference time (e.g., admission time being $t=0$).\par}

We use large language models (LLMs; \texttt{LLaMA\~3.3\~70B}, \texttt{DeepSeek-R1}) to convert PubMed Open Access (PMOA) case reports into \emph{textual time series} (TTS).
A textual time series is defined as an ordered sequence of clinical event–time pairs, where each event is represented as a short free-text span and each timestamp is a numeric value measured in hours relative to a case-specific reference time.
\ChFiveRebuttal{
The reference event is an anchor for the timeline rather than its boundary. For the LLM extraction, hospital admission is assigned $t=0$ when it is described. If admission is not described, the current clinical presentation that frames the case, such as the main diagnosis or treatment, is assigned $t=0$. Events before and after this reference event receive negative and positive timestamps, respectively. Thus, a phrase such as ``three days before admission'' is represented as $-72$ hours, whereas an event on hospital day 2 is represented as $+24$ hours. Depending on the source report, the resulting TTS may cover an acute encounter or a longitudinal course extending over months or years.
}

In this work, a clinical event refers to any health-related mention that can stand alone semantically and is directly relevant to the patient’s course.
This includes symptoms and signs, diagnoses, procedures and diagnostic tests, treatments and medication administrations, major clinical states, and outcomes.
Explicitly stated pertinent negatives and termination events (such as “no shortness of breath” or “antibiotics discontinued”) are also included. 
Demographic attributes that appear in the narrative, such as age and sex are also extracted as clinical events with reference time (t=0) as the timestamp.
We exclude purely contextual text that does not describe the index patient, such as general background statements or literature discussion.


\ChFiveRebuttal{
PMOA-TTS uses source-faithful, span-based events in the tradition of i2b2 clinical concept annotation, while allowing enough surrounding context for an event to remain clinically interpretable, in contrast to span-based temporal annotation methods (e.g., \cite{sun2013evaluating,uzuner20112010}) and i2b2-style guidelines that often limit events to short prepositional phrases. This is necessary because the timelines include symptoms, findings, procedures, treatments, outcomes, pertinent negatives, uncertainty, and intent in addition to diagnoses. For example, retaining ``pain in chest that radiates substernally'' preserves information that would be lost if the event were reduced to a generic chest-pain code; similarly, retaining ``planned'' distinguishes an intended procedure from one that occurred.
}
Conjunctive mentions are separated into individual events when doing so improves clarity.

Natural-language time expressions are normalized into hour units: coarse expressions such as “hospital day 2” are converted using 24-hour increments from the reference point, interval expressions are represented using the start of the interval, and vague temporal phrases are resolved to approximate offsets that are consistent with narrative ordering and surrounding cues.
The goal of this normalization is to produce a usable relative timeline rather than an exact reconstruction of clock time.

The above specifications define the intended representation of textual time series.
The LLM extraction prompt used in our pipeline (Supplementary Material, Section~\ref{ch05:apd:extraction-prompt}) implements this specification by instructing the model to output a structured sequence of \mbox{(clinical event, time)} tuples.
Events are kept in free-text form to retain clinical specificity rather than being normalized to a fixed ontology.

\ChFiveParagraph{Narrative order and temporal order}

We explicitly address the misalignment between narrative structure and chronological reality---where retrospective summaries often reveal outcomes before precursors.

This creates a potential information-leakage risk when retrospective narrative information is treated as available at prediction or decision time.

% !TeX root = ../../../main.tex
\begin{figure}[htbp]
\begin{ChFiveFloatBody}

  \centering

  \subfloat[{Case narrative with highlighted clinical events}]{\begin{minipage}[t]{0.49\textwidth}
    \begin{tcolorbox}
    [
    colback=white,
    colframe=black!30,
    arc=2mm,
    boxrule=0.6pt,
    left=2mm,right=2mm,top=2mm,bottom=2mm
    ]
      \small\raggedright
      A \hl{52-year-old} \hl{woman} was \hl{admitted to the medical ICU} with \hl{chest pain} and \hl{shortness of breath}. She developed \hl{intermittent fever and fatigue} a couple of weeks prior to admission. Eight hours after hospital and ICU admission, \hl{broad-spectrum antibiotics} were \hl{started for suspected infection}. On hospital admission day 2, she developed \hl{pleuritic pain} and \hl{escalating oxygen requirements}. \hl{CT angiography} was performed two days after admission.
      
    \end{tcolorbox}
    
    \label{ch05:fig:narrative_box}
  \end{minipage}}\hfill%
  \subfloat[{Text-ordered event–time tuples.}]{\begin{minipage}[t]{0.49\textwidth}
    \begin{tcolorbox}
    [
    colback=white,
    colframe=black!30,
    arc=2mm,
    boxrule=0.6pt,
    left=2mm,right=2mm,top=2mm,bottom=2mm
    ]
      \small\centering
      \setlength{\tabcolsep}{5pt}
      \begin{adjustbox}{max width=\linewidth}\begin{tabular}{@{} l r @{}}
        \toprule
        \textbf{Event} & \textbf{Time (h)} \\
        \midrule
        52-year-old & 0 \\
        woman & 0 \\
        admitted to the medical ICU & 0 \\
        chest pain & 0 \\
        shortness of breath & 0 \\
        developed intermittent fever & -336 \\
        developed fatigue & -336 \\
        broad-spectrum antibiotics started \\ for suspected infection & 8 \\
        pleuritic pain developed & 24 \\
        escalating oxygen requirements & 24 \\
        CT angiography & 48 \\
        %Management decisions & 48 \\
        \bottomrule
      \end{tabular}\end{adjustbox}
    \end{tcolorbox}
    
    \label{ch05:fig:table_box}
  \end{minipage}}

  
  
  

\end{ChFiveFloatBody}
\caption{Example case report (left) with text-ordered event-time tuples (right).
  }\ChFiveFloatLabel{ch05:fig:annotation_ex}
\end{figure}
